\documentclass[10pt,a4paper]{article}
\usepackage[utf8]{inputenc}
\usepackage[top=2cm,bottom=2cm,left=2cm,right=2cm]{geometry}
\usepackage{amsmath}
\usepackage{amssymb}
\usepackage{array}
\usepackage[dvipsnames,svgnames]{xcolor}
\usepackage{hyperref}

% Anchor persamaan unik meski nomor di-reset per soal.
\newcounter{eqblock}
\newcommand{\reseteq}{%
  \stepcounter{eqblock}%
  \setcounter{equation}{0}%
  \renewcommand{\theHequation}{\arabic{eqblock}.\arabic{equation}}%
}

\newcommand{\dg}[1]{\textcolor{Red}{#1}}

\title{Tugas Minggu 13 : Vektor di Bidang dan Ruang}
\author{Matematika untuk Sistem Informasi}
\date{}

\begin{document}

\maketitle

\begin{enumerate}

%===========================================================
\item Diketahui $P_1(3, -7, 2)$ sebagai titik awal dan $P_2(-2, 5, -4)$ sebagai titik akhir di ruang $\mathbb{R}^3$. Tentukan:
\begin{enumerate}

%-----------------------------------------------------------
\item Komponen vektor $\overrightarrow{P_1P_2}$.

\textbf{Jawab:}
\begin{itemize}
\item \textbf{Langkah 1: Ujung Dikurangi Pangkal}

Komponen vektor diperoleh dari koordinat titik akhir dikurangi koordinat titik awal:
\reseteq
\begin{align}
\overrightarrow{P_1P_2} &= (x_2 - x_1,\; y_2 - y_1,\; z_2 - z_1) \nonumber \\
&= (-2 - 3,\; 5 - (-7),\; -4 - 2) = (-5, 12, -6)
\end{align}
Dalam vektor satuan standar, $\overrightarrow{P_1P_2} = -5\mathbf{i} + 12\mathbf{j} - 6\mathbf{k}$.
% final: (-5, 12, -6)
\end{itemize}

%-----------------------------------------------------------
\item Jarak di antara titik $P_1$ dan $P_2$.

\textbf{Jawab:}
\begin{itemize}
\item \textbf{Langkah 1: Jarak sebagai Norm Vektor Penghubung}

Jarak dua titik sama dengan panjang vektor yang menghubungkannya. Gunakan hasil (1):
\begin{align}
d(P_1, P_2) &= \|\overrightarrow{P_1P_2}\| = \sqrt{(-5)^2 + 12^2 + (-6)^2} \nonumber \\
&= \sqrt{25 + 144 + 36} = \sqrt{205} \approx 14{,}32
\end{align}
Jadi jarak antara $P_1$ dan $P_2$ adalah $\sqrt{205} \approx 14{,}32$ satuan.
% final: \sqrt{205}
\end{itemize}

\end{enumerate}

%===========================================================
\item Diketahui $\mathbf{u} = (-3, 1, 2)$, $\mathbf{v} = (4, 0, 8)$, dan $\mathbf{w} = (6, -1, -4)$ adalah vektor-vektor di $\mathbb{R}^3$. Tentukan:
\begin{enumerate}

%-----------------------------------------------------------
\item $(\mathbf{u} - \mathbf{v}) - \mathbf{w}$

\textbf{Jawab:}
\begin{itemize}
\item \textbf{Langkah 1: Hitung $\mathbf{u} - \mathbf{v}$}

Pengurangan vektor dilakukan komponen demi komponen:
\reseteq
\begin{align}
\mathbf{u} - \mathbf{v} &= (-3 - 4,\; 1 - 0,\; 2 - 8) = (-7, 1, -6)
\end{align}

\item \textbf{Langkah 2: Kurangkan $\mathbf{w}$}

Kurangkan $\mathbf{w}$ dari hasil (1):
\begin{align}
(\mathbf{u} - \mathbf{v}) - \mathbf{w} &= (-7 - 6,\; 1 - (-1),\; -6 - (-4)) = (-13, 2, -2)
\end{align}
Hasilnya berupa \textbf{vektor}.
% final: (-13, 2, -2)
\end{itemize}

%-----------------------------------------------------------
\item $(\|\mathbf{u}\| - \|\mathbf{v}\|) - \|\mathbf{w}\|$

\textbf{Jawab:}
\begin{itemize}
\item \textbf{Langkah 1: Hitung Norm Setiap Vektor}

Kuadratkan setiap komponen, jumlahkan, lalu akarkan:
\reseteq
\begin{align}
\|\mathbf{u}\| &= \sqrt{(-3)^2 + 1^2 + 2^2} = \sqrt{9 + 1 + 4} = \sqrt{14} \\
\|\mathbf{v}\| &= \sqrt{4^2 + 0^2 + 8^2} = \sqrt{16 + 0 + 64} = \sqrt{80} = 4\sqrt{5} \\
\|\mathbf{w}\| &= \sqrt{6^2 + (-1)^2 + (-4)^2} = \sqrt{36 + 1 + 16} = \sqrt{53}
\end{align}

\item \textbf{Langkah 2: Kurangkan Ketiga Bilangan}

Norm adalah bilangan real (skalar), sehingga yang dihitung di sini adalah selisih bilangan biasa:
\begin{align}
(\|\mathbf{u}\| - \|\mathbf{v}\|) - \|\mathbf{w}\| &= \sqrt{14} - 4\sqrt{5} - \sqrt{53} \nonumber \\
&\approx 3{,}742 - 8{,}944 - 7{,}280 \approx -12{,}48
\end{align}
Hasilnya berupa \textbf{skalar} dan boleh negatif karena bukan panjang suatu vektor. Bandingkan dengan panjang hasil (a): $\|(-13, 2, -2)\| = \sqrt{177} \approx 13{,}30$; nilainya berbeda karena norm tidak dapat dikurangkan komponen demi komponen.
% final: \sqrt{14} - 4\sqrt{5} - \sqrt{53}
\end{itemize}

\end{enumerate}

%===========================================================
\item Tentukan apakah $\mathbf{u}$ dan $\mathbf{v}$ membentuk sudut lancip, sudut tumpul, atau saling tegak lurus.
\[
\text{(a) } \mathbf{u} = (1, -5, 4),\; \mathbf{v} = (3, 3, 3) \qquad
\text{(b) } \mathbf{u} = (2, 3),\; \mathbf{v} = (5, -7) \qquad
\text{(c) } \mathbf{u} = (-6, 0, 4),\; \mathbf{v} = (3, 1, 6)
\]

\textbf{Jawab:}
\begin{itemize}
\item \textbf{Langkah 1: Aturan Tanda Hasil Kali Titik}

Dari $\cos\theta = \dfrac{\mathbf{u} \cdot \mathbf{v}}{\|\mathbf{u}\|\,\|\mathbf{v}\|}$ dan penyebut yang selalu positif, jenis sudut cukup ditentukan oleh tanda $\mathbf{u} \cdot \mathbf{v}$ (Bagian~7 Topik~11):
\[
\mathbf{u} \cdot \mathbf{v} > 0 \Rightarrow \text{lancip}, \qquad
\mathbf{u} \cdot \mathbf{v} = 0 \Rightarrow \text{tegak lurus}, \qquad
\mathbf{u} \cdot \mathbf{v} < 0 \Rightarrow \text{tumpul}.
\]

\item \textbf{Langkah 2: Hitung Hasil Kali Titik Setiap Pasangan}

Kalikan komponen yang seletak lalu jumlahkan:
\reseteq
\begin{align}
\text{(a)}\quad \mathbf{u} \cdot \mathbf{v} &= (1)(3) + (-5)(3) + (4)(3) = 3 - 15 + 12 = 0 \\
\text{(b)}\quad \mathbf{u} \cdot \mathbf{v} &= (2)(5) + (3)(-7) = 10 - 21 = -11 \\
\text{(c)}\quad \mathbf{u} \cdot \mathbf{v} &= (-6)(3) + (0)(1) + (4)(6) = -18 + 0 + 24 = 6
\end{align}

\item \textbf{Langkah 3: Simpulkan}

Berdasarkan tanda pada (1)--(3):
\begin{itemize}
\item[(a)] $\mathbf{u} \cdot \mathbf{v} = 0$, sehingga $\mathbf{u}$ dan $\mathbf{v}$ \textbf{saling tegak lurus} ($\theta = 90^\circ$).
% final: \text{saling tegak lurus}
\item[(b)] $\mathbf{u} \cdot \mathbf{v} = -11 < 0$, sehingga $\mathbf{u}$ dan $\mathbf{v}$ membentuk \textbf{sudut tumpul} (vektor di $\mathbb{R}^2$).
% final: \text{sudut tumpul}
\item[(c)] $\mathbf{u} \cdot \mathbf{v} = 6 > 0$, sehingga $\mathbf{u}$ dan $\mathbf{v}$ membentuk \textbf{sudut lancip}.
% final: \text{sudut lancip}
\end{itemize}
\end{itemize}

%===========================================================
\item Tentukan vektor yang tegak lurus terhadap vektor
\[
\mathbf{u} = \begin{bmatrix} -7 \\ 3 \\ 1 \end{bmatrix} \quad \text{dan} \quad
\mathbf{v} = \begin{bmatrix} 2 \\ 0 \\ 4 \end{bmatrix}.
\]

\textbf{Jawab:}
\begin{itemize}
\item \textbf{Langkah 1: Gunakan Hasil Kali Silang}

Hasil kali silang $\mathbf{u} \times \mathbf{v}$ selalu tegak lurus terhadap $\mathbf{u}$ maupun $\mathbf{v}$ (Bagian~8.1 Topik~11). Baris kedua determinan berisi $\mathbf{u}$, baris ketiga berisi $\mathbf{v}$:
\reseteq
\begin{align}
\mathbf{u} \times \mathbf{v} &= \begin{vmatrix} \mathbf{i} & \mathbf{j} & \mathbf{k} \\ -7 & 3 & 1 \\ 2 & 0 & 4 \end{vmatrix}
= \begin{vmatrix} 3 & 1 \\ 0 & 4 \end{vmatrix}\mathbf{i}
- \begin{vmatrix} -7 & 1 \\ 2 & 4 \end{vmatrix}\mathbf{j}
+ \begin{vmatrix} -7 & 3 \\ 2 & 0 \end{vmatrix}\mathbf{k}
\end{align}

\item \textbf{Langkah 2: Hitung Ketiga Minor}

Setiap determinan $2 \times 2$ dihitung dengan $ad - bc$:
\begin{align}
\mathbf{u} \times \mathbf{v} &= (12 - 0)\mathbf{i} - (-28 - 2)\mathbf{j} + (0 - 6)\mathbf{k} \nonumber \\
&= 12\mathbf{i} + 30\mathbf{j} - 6\mathbf{k} = (12, 30, -6)
\end{align}

\item \textbf{Langkah 3: Cek Tegak Lurus}

Hasil kali titik $\mathbf{w} = (12, 30, -6)$ dengan kedua vektor asal harus nol:
\begin{align}
\mathbf{u} \cdot \mathbf{w} &= (-7)(12) + (3)(30) + (1)(-6) = -84 + 90 - 6 = 0 \\
\mathbf{v} \cdot \mathbf{w} &= (2)(12) + (0)(30) + (4)(-6) = 24 + 0 - 24 = 0
\end{align}
Jadi vektor yang tegak lurus terhadap $\mathbf{u}$ dan $\mathbf{v}$ adalah
\[
\mathbf{u} \times \mathbf{v} = (12, 30, -6) = 6\,(2, 5, -1).
\]
Setiap kelipatan $k(2, 5, -1)$ dengan $k \neq 0$ juga tegak lurus terhadap $\mathbf{u}$ dan $\mathbf{v}$, karena mengalikan dengan skalar hanya mengubah panjang (dan mungkin membalik arah), bukan garis kerjanya.
% final: (12, 30, -6)
\end{itemize}

%===========================================================
\item Diketahui vektor
\[
\mathbf{u} = \begin{bmatrix} 2 \\ -1 \\ 3 \end{bmatrix}, \quad
\mathbf{v} = \begin{bmatrix} 4 \\ -1 \\ 2 \end{bmatrix}, \quad \text{dan} \quad
\mathbf{w} = \begin{bmatrix} 2 \\ 6 \\ 7 \end{bmatrix}.
\]
Tentukan hasil kali triple skalar $\mathbf{u} \cdot (\mathbf{v} \times \mathbf{w})$.

\textbf{Jawab:}
\begin{itemize}
\item \textbf{Langkah 1: Determinan dengan Ekspansi Baris Pertama}

Triple skalar sama dengan determinan matriks yang baris-barisnya $\mathbf{u}$, $\mathbf{v}$, $\mathbf{w}$ (Bagian~8.2 Topik~11):
\reseteq
\begin{align}
\mathbf{u} \cdot (\mathbf{v} \times \mathbf{w}) &= \begin{vmatrix} 2 & -1 & 3 \\ 4 & -1 & 2 \\ 2 & 6 & 7 \end{vmatrix}
= 2\begin{vmatrix} -1 & 2 \\ 6 & 7 \end{vmatrix}
- (-1)\begin{vmatrix} 4 & 2 \\ 2 & 7 \end{vmatrix}
+ 3\begin{vmatrix} 4 & -1 \\ 2 & 6 \end{vmatrix} \nonumber \\
&= 2(-7 - 12) + 1(28 - 4) + 3(24 + 2) \nonumber \\
&= -38 + 24 + 78 = 64
\end{align}

\item \textbf{Langkah 2: Cara Alternatif, Silang lalu Titik}

Hitung $\mathbf{v} \times \mathbf{w}$ terlebih dahulu:
\begin{align}
\mathbf{v} \times \mathbf{w} &= \begin{vmatrix} \mathbf{i} & \mathbf{j} & \mathbf{k} \\ 4 & -1 & 2 \\ 2 & 6 & 7 \end{vmatrix}
= (-7 - 12)\mathbf{i} - (28 - 4)\mathbf{j} + (24 + 2)\mathbf{k} = (-19, -24, 26)
\end{align}
Kemudian kalikan titik dengan $\mathbf{u}$:
\begin{align}
\mathbf{u} \cdot (\mathbf{v} \times \mathbf{w}) &= (2)(-19) + (-1)(-24) + (3)(26) = -38 + 24 + 78 = 64
\end{align}
Hasilnya sama dengan (1).

\item \textbf{Langkah 3: Tafsiran Geometris}

Karena $\mathbf{u} \cdot (\mathbf{v} \times \mathbf{w}) = 64 \neq 0$, ketiga vektor \textbf{tidak sebidang}. Nilai mutlaknya adalah volume paralelepipedum yang dibentuk $\mathbf{u}$, $\mathbf{v}$, dan $\mathbf{w}$:
\begin{align}
V &= |\mathbf{u} \cdot (\mathbf{v} \times \mathbf{w})| = |64| = 64 \text{ satuan volume}
\end{align}
Jadi hasil kali triple skalar $\mathbf{u} \cdot (\mathbf{v} \times \mathbf{w}) = 64$.
% final: 64
\end{itemize}

\end{enumerate}

\end{document}
